\section{Fallas comunes y gobernanza}

\subsection{Rasgos comunes de las fallas}

\begin{warningbox}{Falla de generalización durante el entrenamiento}
Aunque obtenga una reward alta en el entorno de entrenamiento, el agent puede fallar ante un estado inicial ligeramente modificado. La causa suele ser que el replay buffer depende en exceso de un pequeño grupo de transiciones con reward alta y no logra cubrir el espacio de estados. Esto puede mitigarse con un entropy bonus o con replay multitarea.
\end{warningbox}

\begin{knowledgebox}{Aspectos de gobernanza}
\begin{itemize}
    \item Policy drift: hay que conservar muestras de la política dev como referencia dentro de la ventana.
    \item Reward hacking: en el aprendizaje por refuerzo, la reward shaping exige cautela; no se debe modificar la reward de forma simple de modo que el agent encuentre un shortcut.
    \item Transparencia de los datos: registrar el slice de datos usado en cada actualización y cotejar periódicamente los registros de comportamiento con las áreas de negocio.
\end{itemize}
\end{knowledgebox}

\subsection{Mecanismos de intervención}

\begin{importantbox}{Safe Interruptibility e intervención humana}
Durante el entrenamiento, mantener la interruptibility significa que, cuando una persona detecta una anomalía, puede detener o reiniciar el agent de forma segura sin provocar una explosión de la recompensa. En la práctica se configura un \texttt{kill switch} y, al mismo tiempo, se registra el último TD error para ayudar a localizar la causa.
\end{importantbox}

\begin{warningbox}{Fragmentación del entrenamiento por intervención excesiva}
Interrumpir el entrenamiento con frecuencia hace que los datos del replay buffer provengan de políticas distintas, lo que rompe el supuesto off-policy. Se recomienda ejecutar un warm-start offline training después de cada intervención, para que la nueva política tenga tiempo suficiente de converger.
\end{warningbox}

\subsection{Resumen de la sección}

En el plano de la gobernanza hay que equilibrar el entrenamiento automático con la supervisión humana; los mecanismos de intervención deben centrarse en las guardrails y evitar la fragmentación por entrenamientos repetidos.
